\section{Herramientas y/o lenguajes de programación}

Dado que la presente etapa del proyecto ``CoConsorcio'' se centró en el relevamiento, la especificación de requisitos y la simulación de la gestión bajo un marco de trabajo ágil, no se emplearon lenguajes de programación ni entornos de desarrollo de código. En su lugar, el equipo hizo uso exclusivo de herramientas de software orientadas a la administración de proyectos, la redacción colaborativa y la comunicación.  

A continuación, se describen brevemente las herramientas utilizadas:

\begin{itemize}
    \item \textbf{Google Docs:} Procesador de textos basado en la nube utilizado para la redacción colaborativa en tiempo real. Fue la herramienta principal para estructurar, revisar y unificar toda la documentación académica del proyecto, incluyendo la Especificación de Requisitos de Software (ERS) y la transcripción de entrevistas.  
    \item \textbf{Trello:} Software de administración de proyectos basado en la metodología de tableros visuales Kanban. Se utilizó para simular y gestionar el ciclo de vida del marco de trabajo Scrum, permitiendo al equipo administrar el \textit{Product Backlog}, planificar los Sprints y realizar el seguimiento del estado de las Historias de Usuario.  
    \item \textbf{Gemini (Inteligencia Artificial):} Asistente basado en modelos de lenguaje de gran tamaño (LLM) empleado como herramienta de apoyo metodológico. Se utilizó para estructurar ideas, refinar la redacción técnica de los requerimientos bajo el estándar IEEE y generar escenarios de simulación (datos de prueba para entrevistas e historias de usuario).  
    \item \textbf{Discord:} Plataforma de comunicación mediante canales de voz y texto. Se empleó para llevar a cabo las reuniones sincrónicas del equipo, especialmente para coordinar las tareas, simular las reuniones ágiles (\textit{Daily Standups}) y compartir pantalla durante la revisión de los documentos.  
    \item \textbf{WhatsApp:} Aplicación de mensajería instantánea utilizada como canal secundario para la comunicación asincrónica y rápida del equipo. Facilitó la organización de horarios, avisos urgentes y el intercambio de enlaces directos a los recursos del proyecto.  
\end{itemize}
